\section{问题重述}

\subsection{问题背景}

无创产前检测（Non-invasive Prenatal Test, NIPT）通过采集母体外周血中的胎儿游离 DNA 片段\cite{ref1}，分析胎儿染色体是否存在异常，可在孕早期筛查唐氏综合征（21 三体）、爱德华氏综合征（18 三体）和帕陶氏综合征（13 三体）\cite{ref5}。NIPT 的准确性主要由胎儿性染色体浓度决定：对于男胎，若 Y 染色体游离 DNA 片段比例（简称 Y 染色体浓度）达到或超过 4\%，且女胎 X 染色体浓度无异常，则认为检测结果基本准确，否则难以保证准确性要求。同时，实际中应尽早发现不健康的胎儿，否则会带来治疗窗口期缩短的风险，早期发现（12 周以内）风险较低，中期发现（13--27 周）风险高，晚期发现（28 周以后）风险极高。

实践表明，男胎 Y 染色体浓度与孕妇孕周数及身体质量指数（BMI）紧密相关。由于每个孕妇的年龄、BMI、孕情等存在个体差异，对所有孕妇采用统一的检测时点进行 NIPT 会显著影响其准确性。因此，依据 BMI 等指标对孕妇进行合理分组、确定各群组的最佳 NIPT 时点，是降低潜在风险的关键。

\subsection{需要解决的问题}

\textbf{问题一：} 分析胎儿 Y 染色体浓度与孕妇孕周数和 BMI 等指标的相关特性，建立相应的关系模型，并检验其显著性。

\textbf{问题二：} 男胎孕妇的 BMI 是影响 Y 染色体浓度最早达标时间（浓度达到或超过 4\% 的最早时间）的主要因素。对男胎孕妇的 BMI 进行合理分组，给出每组的 BMI 区间和最佳 NIPT 时点，使潜在风险最小，并分析检测误差对结果的影响。

\textbf{问题三：} Y 染色体浓度达标时间受身高、体重、年龄等多种因素影响。综合考虑这些因素、检测误差和达标比例，根据 BMI 给出合理分组及每组的最佳 NIPT 时点，使潜在风险最小，并分析检测误差的影响。

\textbf{问题四：} 孕妇和女胎都不携带 Y 染色体，需判定女胎是否异常。以女胎孕妇的 21 号、18 号和 13 号染色体非整倍体（AB 列）为判定结果，综合 X 染色体及上述染色体的 Z 值、GC 含量、读段数及相关比例、BMI 等因素，给出女胎异常的判定方法。
